\section*{A two-state finite-coupling calculation}
\addcontentsline{toc}{section}{A two-state finite-coupling calculation}

The difference between tangent geometry and an exact pair law is already visible for two uniform
binary variables.  Let
\begin{equation}
p_i=p_j=(1/2,1/2)^{\mathsf T}
\end{equation}
and choose the centered score
\begin{equation}
S_\theta
=\theta
\begin{pmatrix}
1&-1\\
-1&1
\end{pmatrix}.
\label{eq:book-binary-score}
\end{equation}
Because symmetry gives equal row and column sums after exponentiation, exact matrix scaling can be
written in closed form:
\begin{equation}
Q_\theta
=\frac{1}{4\cosh\theta}
\begin{pmatrix}
e^\theta&e^{-\theta}\\
e^{-\theta}&e^\theta
\end{pmatrix}
=\frac14
\begin{pmatrix}
1+t&1-t\\
1-t&1+t
\end{pmatrix},
\qquad t=\tanh\theta.
\label{eq:book-binary-scaled-law}
\end{equation}
Every row and column sums exactly to $1/2$ for every finite $\theta$.

Near independence, $t=\theta+O(\theta^3)$ and
\begin{equation}
Q_\theta-p_ip_j^{\mathsf T}
=\frac{\theta}{4}
\begin{pmatrix}
1&-1\\
-1&1
\end{pmatrix}
+O(\theta^3),
\label{eq:book-binary-tangent}
\end{equation}
which is exactly the tangent direction predicted by
Eq.~\ref{eq:sinkhorn-linearization}.  At finite strength the linear coordinate $\theta$ is replaced
in the probability table by the bounded coordinate $\tanh\theta$.  The transport polytope bends
and eventually saturates as probability concentrates on the two matching states.

The mutual information is also exact:
\begin{equation}
I(\theta)
=\frac{1+t}{2}\log(1+t)
+\frac{1-t}{2}\log(1-t).
\label{eq:book-binary-mutual-information}
\end{equation}
For weak coupling, $I(\theta)=\theta^2/2+O(\theta^4)$, agreeing with the quadratic tangent norm.
At strong coupling, the nonlinear probability geometry matters and the quadratic approximation no
longer describes the finite law.

Finally, add arbitrary row and column scores to $S_\theta$.  They change the unscaled exponential
matrix, but exact marginal scaling absorbs them and returns the same $Q_\theta$.  The nonlinear
probability model therefore respects precisely the same gauge quotient as the linear operator,
while encoding finite displacement on that quotient rather than only its tangent direction.
